\documentclass[12pt,a4paper]{article}

% ─── Paquetes ────────────────────────────────────────────────────────────────
\usepackage[utf8]{inputenc}
\usepackage[T1]{fontenc}
\usepackage[spanish]{babel}
\usepackage{geometry}
\usepackage{booktabs}
\usepackage{longtable}
\usepackage{array}
\usepackage{xcolor}
\usepackage{hyperref}
\usepackage{enumitem}
\usepackage{titlesec}
\usepackage{fancyhdr}
\usepackage{graphicx}
\usepackage[activate=false]{microtype}
\usepackage{parskip}

% ─── Geometría ───────────────────────────────────────────────────────────────
\geometry{
  top=2.5cm, bottom=2.5cm,
  left=3cm, right=3cm
}

% ─── Colores ─────────────────────────────────────────────────────────────────
\definecolor{primary}{HTML}{1A5276}
\definecolor{secondary}{HTML}{2E86C1}
\definecolor{accent}{HTML}{F39C12}
\definecolor{lightgray}{HTML}{F2F3F4}
\definecolor{darkgray}{HTML}{424949}

% ─── Tipografía de secciones ─────────────────────────────────────────────────
\titleformat{\section}
  {\large\bfseries\color{primary}}
  {\thesection.}{0.5em}{}[\titlerule]

\titleformat{\subsection}
  {\normalsize\bfseries\color{secondary}}
  {\thesubsection.}{0.5em}{}

% ─── Encabezado y pie ────────────────────────────────────────────────────────
\pagestyle{fancy}
\fancyhf{}
\rhead{\textcolor{darkgray}{\small Red Yerba Buena --- CCC Grant 2026}}
\lhead{\textcolor{darkgray}{\small Colectivo Nodo Norte / cheLA / Fundación exACTa}}
\rfoot{\textcolor{darkgray}{\small Página \thepage}}
\renewcommand{\headrulewidth}{0.4pt}

% ─── Hiperlinks ──────────────────────────────────────────────────────────────
\hypersetup{
  colorlinks=true,
  linkcolor=primary,
  urlcolor=secondary,
  pdfauthor={Patricia Vanesa Quiroz},
  pdftitle={Red Yerba Buena — Propuesta CCC 2026}
}

% ═══════════════════════════════════════════════════════════════════════════════
\begin{document}
% ═══════════════════════════════════════════════════════════════════════════════

% ─── Portada ─────────────────────────────────────────────────────────────────
\begin{titlepage}
  \centering
  \vspace*{2cm}

  {\color{primary}\rule{\linewidth}{2pt}}
  \vspace{0.5cm}

  {\Huge\bfseries\color{primary} Red Yerba Buena}\\[0.4cm]
  {\large\color{secondary} Ecosistema de Autonomía Territorial\\y Resonancia Aplicada}\\[0.3cm]
  {\normalsize\color{darkgray} Quebrada de Humahuaca --- Jujuy --- Argentina}

  \vspace{0.5cm}
  {\color{primary}\rule{\linewidth}{2pt}}

  \vspace{2cm}

  \begin{tabular}{ll}
    \textbf{Convocatoria:}   & Internet Society Foundation --- CCC Grant 2026 \\[4pt]
    \textbf{Track:}          & Catalyst Track (hasta USD 50.000) \\[4pt]
    \textbf{Monto solicitado:} & USD 50.000 \\[4pt]
    \textbf{Duración:}       & 12 meses \\[4pt]
    \textbf{Organización:}   & Colectivo Nodo Norte \\[4pt]
    \textbf{Entidad fiscal:} & Fundación exACTa / cheLA \\[4pt]
    \textbf{Responsable:}    & Patricia Vanesa Quiroz \\[4pt]
    \textbf{Contacto:}       & \href{mailto:vaneq@colectivonodonorte.ar}{vaneq@colectivonodonorte.ar} \\[4pt]
    \textbf{Fecha:}          & Mayo 2026 \\
  \end{tabular}

  \vfill

  {\small\color{darkgray}
    Colectivo Nodo Norte opera como nodo BioCulturas de cheLA\\
    (Centro Hipermediático Experimental Latinoamericano)\\
    bajo la personería jurídica de \textbf{Fundación exACTa}\\
    en colaboración con \textbf{UCLA REMAP}
  }

  \vspace{1cm}
  {\color{primary}\rule{\linewidth}{1pt}}
\end{titlepage}

% ─── Índice ──────────────────────────────────────────────────────────────────
\tableofcontents
\newpage

% ═══════════════════════════════════════════════════════════════════════════════
\section{Resumen Ejecutivo}
% ═══════════════════════════════════════════════════════════════════════════════

\textbf{Red Yerba Buena} es una red de conectividad \textit{mesh} LoRa (tecnología Meshtastic) diseñada para siete comunidades indígenas de la Quebrada de Humahuaca, provincia de Jujuy, Argentina, donde la telefonía móvil es inexistente debido a la compleja orografía andina.

La solución opera sin infraestructura corporativa, sin costos de mantenimiento recurrentes y sin dependencia de conectividad satelital. Será gestionada y mantenida por mujeres rurales de cada paraje, garantizando la autonomía operativa de la red.

El proyecto es implementado por \textbf{Colectivo Nodo Norte} bajo el marco institucional de \textbf{Fundación exACTa / cheLA}, con más de 20 años de experiencia en tecnología comunitaria en América Latina en colaboración con \textbf{UCLA REMAP}.

Esta solicitud financia la \textbf{Etapa 1} de un proyecto de USD 250.000 en tres etapas. La Etapa 1 despliega la infraestructura base en las siete comunidades y establece la capacidad institucional y comunitaria necesaria para escalar.

% ═══════════════════════════════════════════════════════════════════════════════
\section{Contexto y Problema}
% ═══════════════════════════════════════════════════════════════════════════════

\subsection{El territorio}

La Quebrada de Humahuaca es un valle andino de 155 km de longitud declarado \textbf{Patrimonio de la Humanidad por la UNESCO en 2003}. Sus comunidades originarias, reconocidas constitucionalmente desde 1994 bajo el Artículo 75 inciso 17 de la Constitución Nacional Argentina, habitan parajes de alta montaña a altitudes que oscilan entre los 2.000 y los 3.800 metros sobre el nivel del mar.

\subsection{La brecha de conectividad}

La orografía andina imposibilita el despliegue de infraestructura de telefonía móvil convencional en los parajes rurales de la Quebrada. Las familias que habitan estos territorios no tienen acceso a comunicación digital de ningún tipo, lo que genera:

\begin{itemize}
  \item Imposibilidad de comunicar emergencias sanitarias o climáticas
  \item Aislamiento respecto de servicios públicos y educativos
  \item Exclusión de oportunidades económicas digitales
  \item Éxodo de jóvenes hacia centros urbanos por falta de horizontes locales
\end{itemize}

\subsection{Conectividad de emergencia: Starlink y la escuela como excepción}

En los escasos casos en que algún emprendimiento productivo o turístico tiene acceso a internet, este se obtiene exclusivamente mediante la contratación individual de \textbf{Starlink} —a un costo mensual prohibitivo para la economía rural andina— o mediante el desplazamiento físico hasta la escuela más cercana, que en algunos parajes cuenta con conexión \textit{wifi}. En ambos casos se trata de un acceso fragmentado, costoso e impracticable para la comunicación cotidiana con clientes y proveedores.

Esta situación tiene consecuencias directas sobre la viabilidad de los emprendimientos familiares y comunitarios: la imposibilidad de mantener comunicación fluida con turistas, agencias y mercados impide la consolidación de cualquier actividad económica sostenible en el territorio.

Meshtastic resuelve este problema de dos maneras complementarias: por un lado, provee comunicación autónoma entre parajes \textit{sin necesidad de internet}; por otro, cuando existe un nodo con acceso a internet (Starlink o escuela), ese nodo funciona como \textbf{gateway MQTT} y extiende la conectividad hacia aplicaciones externas como \textbf{WhatsApp} (vía protocolo MQTT-bridge), permitiendo que toda la red se beneficie de ese acceso puntual sin que cada hogar deba costear la conexión de forma individual.

\subsection{El contexto humano}

Desde 2003 la UNESCO reconoce el valor patrimonial de este territorio. Sin embargo, las familias jóvenes se ven forzadas al éxodo urbano por falta de caminos, energía, transporte y conectividad. Este proyecto existe para que los jóvenes tengan un lugar adonde volver, y para que formar una familia en el territorio sea un horizonte posible.

\subsection{Las comunidades beneficiarias}

El proyecto interviene en siete comunidades distribuidas a lo largo de la Quebrada de Humahuaca:

\begin{center}
\begin{tabular}{clc}
\toprule
\textbf{N°} & \textbf{Comunidad} & \textbf{Provincia} \\
\midrule
1 & Humahuaca       & Jujuy \\
2 & Chorrillos      & Jujuy \\
3 & Hornaditas      & Jujuy \\
4 & Negra Muerta    & Jujuy \\
5 & Iturbe          & Jujuy \\
6 & Azul Pampa      & Jujuy \\
7 & Tres Cruces     & Jujuy \\
\bottomrule
\end{tabular}
\end{center}

% ═══════════════════════════════════════════════════════════════════════════════
\section{La Solución}
% ═══════════════════════════════════════════════════════════════════════════════

\subsection{Tecnología Meshtastic}

Meshtastic es un protocolo de comunicación de código abierto que opera sobre radiofrecuencia LoRa (\textit{Long Range}). Permite transmitir mensajes de texto, coordenadas GPS y alertas entre dispositivos sin necesidad de internet, telefonía móvil ni infraestructura centralizada.

Cada nodo amplifica y retransmite la señal de sus vecinos, creando una red mallada (\textit{mesh}) que se extiende a lo largo de la cadena montañosa. La señal LoRa puede alcanzar distancias de 10 a 30 km en condiciones de línea de vista, lo que la hace especialmente adecuada para la orografía andina.

\subsection{Puente hacia WhatsApp mediante MQTT}

Una de las capacidades más relevantes de Meshtastic para el contexto productivo y turístico de las comunidades es su compatibilidad con el protocolo \textbf{MQTT} (\textit{Message Queuing Telemetry Transport}). Cuando uno o más nodos de la red tienen acceso a internet —ya sea por Starlink, conexión escolar u otra fuente puntual— ese nodo actúa como \textbf{gateway} y publica los mensajes de la red \textit{mesh} en un servidor MQTT. Desde allí, mediante integraciones disponibles, los mensajes pueden entregarse en \textbf{grupos de WhatsApp} accesibles desde cualquier celular con señal convencional.

Esto tiene implicancias concretas para los emprendimientos comunitarios:

\begin{itemize}
  \item Un turista puede consultar disponibilidad de alojamiento desde su teléfono con señal, y la consulta llega al poblador en el cerro a través de la red \textit{mesh}.
  \item Un emprendimiento puede recibir pedidos o coordinar con proveedores sin moverse del paraje.
  \item La red entera se beneficia de cada punto de acceso a internet existente, sin que cada hogar deba costear conectividad individual.
\end{itemize}

La configuración MQTT es parte del proceso de instalación y capacitación previsto en el plan de trabajo.

\subsection{Arquitectura de la red}

Cada comunidad contará con dos tipos de dispositivos:

\begin{description}
  \item[Terminales personales] Dispositivos compactos (USD 30 c/u) que los pobladores llevan consigo y conectan a su teléfono celular mediante Bluetooth. Permiten enviar y recibir mensajes sin señal de celular ni internet.

  \item[Nodos fijos en altura] Equipos con carcasa de protección IP67, panel solar y batería (USD 90 c/u) instalados en puntos elevados de los cerros. Garantizan la cobertura y el alcance de la red entre parajes. Los nodos ubicados en puntos con acceso a internet (escuelas, parajes con Starlink) funcionan además como gateways MQTT.
\end{description}

\subsection{Infraestructura de movilidad autónoma}

El trabajo de campo en la Quebrada requiere movilidad en un territorio donde el combustible fósil es escaso y costoso. El proyecto incorpora una \textbf{Ford F-100 adaptada con tecnología de gasificación de biomasa} según los planos de código abierto de \href{https://autoabasura.com}{autoabasura.com}, desarrollados por el Ing. Edmundo Ramos.

Este sistema convierte residuos orgánicos secos en gas combustible (\textit{gasura}), permitiendo operar el vehículo sin gasolina. Los gases de escape presentan una pureza superior al aire atmosférico, con aporte de oxígeno al entorno. La capacitación e instalación se realizará directamente con el Ing. Ramos en su taller de Córdoba.

\subsection{Gestión comunitaria por mujeres}

La red será operada y mantenida por mujeres rurales de cada paraje. El diseño del sistema prioriza la simplicidad operativa: los nodos fijos son autónomos (panel solar, sin intervención diaria) y las terminales personales funcionan como cualquier dispositivo Bluetooth.

% ═══════════════════════════════════════════════════════════════════════════════
\section{Participación Comunitaria}
% ═══════════════════════════════════════════════════════════════════════════════

Colectivo Nodo Norte cuenta con vínculos comunitarios profundos y de larga data en el territorio. La coordinadora del proyecto, Patricia Vanesa Quiroz, fue \textbf{docente de Tecnología en la Tecnicatura de Desarrollo Indígena} del \textbf{COAJ} (Centro de Organizaciones Aborígenes de Jujuy), institución que nuclea a las principales organizaciones indígenas de la provincia. Fue también docente de diseño gráfico y video en la \textbf{Tecnicatura de Comunicación Social del IES de Humahuaca}, y acompañó proyectos productivos comunitarios ---entre ellos la producción de harina de habas en \textbf{Azul Pampa}--- en el marco de la carrera de Técnico Universitario en Valor Agregado en Proyecto Productivo.

Como resultado de esta trayectoria docente, muchos de los \textbf{referentes comunitarios actuales de los siete parajes son ex alumnos directos} de la coordinadora, lo que garantiza relaciones de confianza preexistentes, conocimiento mutuo del territorio y comprensión compartida de las necesidades tecnológicas de las comunidades.

El proyecto está transitando el proceso formal de \textbf{consulta libre, previa e informada} con cada asamblea comunitaria, respetando el Convenio 169 de la OIT incorporado a la Constitución Argentina (Art. 75 inc. 17). Este proceso es inherentemente participativo y no puede acelerarse sin comprometer la soberanía comunitaria. Como primer resultado concreto, \textbf{en mayo de 2026 se firma el convenio entre Fundación exACTa y la comunidad de Azul Pampa}, estableciendo el marco formal de colaboración para el despliegue de la red en ese paraje.

La metodología \textbf{A.B.R.A.} (Acciones Biosistémicas de Resonancia Aplicada) que guía el trabajo del Colectivo prioriza el co-diseño con la comunidad antes de cualquier despliegue técnico. La red será co-diseñada, implementada y mantenida por las propias comunidades, garantizando la sostenibilidad del proyecto más allá del período de financiamiento.

% ═══════════════════════════════════════════════════════════════════════════════
\section{Marco Institucional}
% ═══════════════════════════════════════════════════════════════════════════════

\subsection{Colectivo Nodo Norte}

Laboratorio de innovación social y soberanía tecnológica radicado en la Quebrada de Humahuaca. Opera bajo el marco de acción A.B.R.A. (Acciones Biosistémicas de Resonancia Aplicada), un enfoque integral que busca el desarrollo humano consciente mediante el equilibrio entre el ser y su entorno. Se dedica al diseño e implementación de Tecnologías Situadas para revertir el vaciamiento de los territorios rurales, integrando saberes ancestrales con herramientas contemporáneas.

\subsection{Fundación exACTa / cheLA}

\textbf{Fundación exACTa} es una organización sin fines de lucro con sede en Buenos Aires, Argentina, fundada en 2003 en colaboración con \textbf{UCLA REMAP} y el UCLA Program on Digital Cultures. A través de \textbf{cheLA} (Centro Hipermediático Experimental Latinoamericano), desarrolla investigación y práctica en la intersección del arte, la ciencia y la tecnología con impacto en comunidades latinoamericanas.

cheLA se organiza en nodos temáticos de investigación y creación colectiva. Colectivo Nodo Norte actúa como \textbf{nodo BioCulturas} de esta red, bajo cuyo marco se presenta este proyecto.

Entre los proyectos más relevantes de la organización para este programa se destacan: \textbf{EduSol}, programa de educación con Software Libre para escuelas del barrio de Parque Patricios (donde se ubica cheLA), desarrollado por Patricia Vanesa Quiroz, que formó a docentes y estudiantes en herramientas digitales de acceso abierto. \textbf{sur@norte@sur}, programa de investigación financiado por Latin American Digital Cultures de UCLA, del que emergió la investigación de Fabian Wagmister en \textit{Computación Cívica Comunitaria} (\textit{Cultural Civic Computing}), marco conceptual directo del presente proyecto. \textbf{Experiencia Iberá} (2016), primer proyecto en Argentina en usar una red de sensores audiovisuales para una instalación interactiva comunitaria, diseñada con UCLA REMAP y exhibida en Tecnópolis ante decenas de miles de visitantes. \textbf{Las Cuencas Cuentan} (2024--2026), residencia de investigación y creación en ambientes hídricos en territorio, con apoyo de Metabolic Studio y Fundación Williams.

\subsection{Responsable del proyecto}

\textbf{Patricia Vanesa Quiroz} fue coordinadora de cheLA entre 2004 y 2007, período en que se consolidó la alianza institucional entre Fundación exACTa y UCLA REMAP, mediante proyectos de Computación Cultural Comunitaria. Durante ese período desarrolló \textbf{EduSol}, programa de educación con Software Libre para escuelas de Parque Patricios, e integró el equipo del programa de investigación \textbf{sur@norte@sur} (Latin American Digital Cultures, UCLA). Su trayectoria de 20 años en la intersección de tecnología comunitaria, cultura y territorio en América Latina es la base operativa de este proyecto.

\subsection{Asesor académico}

\textbf{Fabian Wagmister}, Director General de cheLA y fundador de UCLA REMAP, actúa como asesor estratégico del proyecto. Su trabajo en \textit{Cultural Civic Computing} y su colaboración activa con comunidades latinoamericanas proveen el marco conceptual y académico de la iniciativa.

% ═══════════════════════════════════════════════════════════════════════════════
\section{Plan de Trabajo}
% ═══════════════════════════════════════════════════════════════════════════════

\subsection{Etapa 1 — Esta solicitud (Meses 1--12)}

\begin{description}
  \item[Meses 1--2: Preparación institucional y técnica]
    Adquisición del vehículo. Viaje a Córdoba con el instalador para capacitación e instalación del gasificador con el Ing. Edmundo Ramos. Compra de hardware Meshtastic (terminales y nodos).

  \item[Meses 2--3: Capacitación técnica]
    Visita del técnico de Meshtastic Argentina para capacitación en instalación, configuración y mantenimiento de nodos, incluyendo la configuración de gateways MQTT para los puntos con acceso a internet. Participan la coordinadora, el instalador y referentes comunitarios.

  \item[Meses 3--10: Despliegue por comunidades]
    Instalación progresiva en las siete comunidades. Dos fines de semana por mes de trabajo de campo. Proceso paralelo de consulta comunitaria y co-diseño de protocolos de uso. Relevamiento de puntos con Starlink o wifi escolar para configuración de gateways MQTT.

  \item[Meses 6--12: Capacitación a mujeres operadoras]
    Formación de referentes femeninas en cada paraje para la gestión autónoma de la red. Monitoreo de uso y ajuste de protocolos.

  \item[Meses 6--12: Diseño y gestión de proyectos productivos y turísticos]
    La coordinadora del proyecto trabaja de forma paralela al despliegue técnico en el \textbf{diseño, gestión y hosting de proyectos productivos y turísticos} familiares y comunitarios: identidad digital, páginas web, canales de comunicación con clientes y proveedores (incluyendo integración WhatsApp--MQTT), y sistemas básicos de reserva y comercialización. Este componente forma parte integral de los honorarios de coordinación, ya que la conectividad cobra sentido únicamente si existe una oferta comunitaria organizada capaz de aprovecharla.

  \item[Meses 10--12: Documentación y evaluación]
    Sistematización de aprendizajes, métricas de uso, informe final. Preparación de la Etapa 2 (Scaling Track).
\end{description}

\subsection{Etapas futuras (referencia)}

\begin{description}
  \item[Etapa 2 --- USD 100.000] Extensión a comunidades adicionales de la Quebrada, Puna y Valles jujeños y consolidación institucional del Colectivo en Jujuy.
  \item[Etapa 3 --- USD 100.000] Replicación del modelo en otras regiones andinas de Argentina.
\end{description}

% ═══════════════════════════════════════════════════════════════════════════════
\section{Presupuesto}
% ═══════════════════════════════════════════════════════════════════════════════

\subsection{Resumen por categoría}

\begin{center}
\begin{longtable}{p{8cm}rr}
\toprule
\textbf{Concepto} & \textbf{Detalle} & \textbf{USD} \\
\midrule
\endfirsthead
\toprule
\textbf{Concepto} & \textbf{Detalle} & \textbf{USD} \\
\midrule
\endhead
\midrule
\multicolumn{3}{r}{\small continúa en la página siguiente}\\
\endfoot
\bottomrule
\endlastfoot

\multicolumn{3}{l}{\textbf{Hardware de red}} \\
Terminales personales Meshtastic
  & 70 unidades × USD 30 & 2.100 \\
Nodos fijos en cerros (panel solar + protección IP67)
  & 28 unidades × USD 90 & 2.520 \\
\midrule

\multicolumn{3}{l}{\textbf{Herramental}} \\
Equipamiento para instalación en altura
  & Taladro, arnés, multímetro, etc. & 800 \\
\midrule

\multicolumn{3}{l}{\textbf{Capacitación técnica}} \\
Honorarios técnico Meshtastic Argentina
  & 1 persona × 1 semana & 1.000 \\
Pasaje aéreo BUE--JUJ--BUE
  & 1 pasaje & 200 \\
Viáticos 7 días
  & Alimentación y movilidad local & 140 \\
\midrule

\multicolumn{3}{l}{\textbf{Movilidad territorial}} \\
Ford F-100 segunda mano
  & Vehículo de campo & 6.000 \\
Capacitación e instalación de gasificador
  & Ing. Edmundo Ramos, Córdoba & 650 \\
Traslados y viáticos (coordinadora + instalador)
  & Viaje a Córdoba & 200 \\
\midrule

\multicolumn{3}{l}{\textbf{Recursos humanos}} \\
Coordinación del proyecto (12 meses)
  & Patricia Vanesa Quiroz & 18.000 \\
\multicolumn{3}{l}{\small\color{darkgray}\itshape
  \quad Incluye: dirección técnica y comunitaria · diseño, gestión} \\
\multicolumn{3}{l}{\small\color{darkgray}\itshape
  \quad y hosting de proyectos productivos y turísticos · integración} \\
\multicolumn{3}{l}{\small\color{darkgray}\itshape
  \quad WhatsApp--MQTT · identidad digital de emprendimientos} \\
Instalador en campo (24 fines de semana + viáticos)
  & Técnico local & 8.640 \\
\midrule

\multicolumn{3}{l}{\textbf{Administración}} \\
Gestión contable y Sponsor Fiscal
  & Fundación exACTa & 5.000 \\
\midrule

\multicolumn{3}{l}{\textbf{Imprevistos}} \\
Reserva de contingencia (8,5\%)
  & Trabajo en alta montaña & 4.250 \\
\midrule[\heavyrulewidth]

\multicolumn{2}{l}{\textbf{TOTAL}} & \textbf{50.000} \\

\end{longtable}
\end{center}

\subsection{Notas al presupuesto}

\begin{itemize}
  \item El vehículo queda en propiedad de la organización y es reutilizado en las Etapas 2 y 3.
  \item Los planos del gasificador son de código abierto y descarga gratuita (\href{https://autoabasura.com}{autoabasura.com}).
  \item El hospedaje durante la capacitación en campo es provisto por la red comunitaria sin costo.
  \item Los honorarios del instalador incluyen viáticos de traslado entre los siete parajes.
  \item Los honorarios de coordinación incluyen el diseño, gestión y hosting de proyectos productivos y turísticos comunitarios, la configuración del bridge WhatsApp--MQTT y la identidad digital de los emprendimientos. Este trabajo es inseparable de la infraestructura técnica: la red cobra valor cuando existe una oferta organizada capaz de aprovecharla.
  \item El porcentaje de imprevistos refleja la complejidad logística del trabajo en alta montaña andina.
\end{itemize}

% ═══════════════════════════════════════════════════════════════════════════════
\section{Sostenibilidad}
% ═══════════════════════════════════════════════════════════════════════════════

La sostenibilidad del proyecto descansa en tres pilares:

\begin{description}
  \item[Autonomía técnica] Los nodos Meshtastic son alimentados por energía solar, no requieren mantenimiento diario y operan indefinidamente sin costos recurrentes. Las mujeres operadoras capacitadas pueden resolver las incidencias más comunes sin asistencia externa.

  \item[Autonomía energética] El vehículo adaptado con gasificador opera con biomasa local (residuos orgánicos secos disponibles en todos los parajes), eliminando la dependencia de combustible fósil para las operaciones de campo.

  \item[Base institucional] La trayectoria de Fundación exACTa/cheLA garantiza la continuidad institucional del proyecto más allá del período de financiamiento. El modelo de nodo BioCulturas permite integrar nuevas fuentes de financiamiento (Etapas 2 y 3) sobre la base construida.

  \item[Autogeneración de ingresos] La red habilita un modelo de generación de ingresos propios que financia su mantenimiento a largo plazo. Colectivo Nodo Norte desarrolla y gestiona \textbf{portales web e identidad digital} para los emprendimientos turísticos y productivos de las comunidades, canalizando la comunicación con turistas y consumidores externos —incluyendo la integración con WhatsApp vía MQTT— hacia reservas y ventas concretas. A cambio, el Colectivo percibe una \textbf{comisión por cada reserva o venta concretada}, sin costo fijo para el emprendimiento mientras no haya resultado. Este modelo alinea los incentivos del Colectivo con el éxito económico de las familias: a mayor actividad turística y productiva en el territorio, mayor la capacidad de autofinanciamiento de la red y del equipo que la sostiene. Con un número suficiente de emprendimientos activos, este flujo cubre los costos operativos de la red sin depender de financiamiento externo.
\end{description}

% ═══════════════════════════════════════════════════════════════════════════════
\section{Impacto Esperado}
% ═══════════════════════════════════════════════════════════════════════════════

Al finalizar la Etapa 1:

\begin{itemize}
  \item \textbf{700 personas} con acceso a comunicación digital en 7 comunidades previamente aisladas
  \item \textbf{70 terminales personales} distribuidas, con prioridad a mujeres y familias con menores
  \item \textbf{28 nodos fijos} instalados en puntos estratégicos de los cerros, incluyendo gateways MQTT en puntos con acceso a internet
  \item \textbf{7 mujeres operadoras} capacitadas para la gestión autónoma de la red
  \item \textbf{Proyectos productivos y turísticos} de familias y comunidades con presencia digital, canales de comunicación activos y acceso a WhatsApp desde la red \textit{mesh}
  \item \textbf{1 vehículo de campo} operativo con movilidad autónoma de combustible fósil
  \item \textbf{Documentación completa} del modelo replicable para otras comunidades andinas
\end{itemize}

% ═══════════════════════════════════════════════════════════════════════════════
\section{Equipo}
% ═══════════════════════════════════════════════════════════════════════════════

\begin{description}
  \item[Patricia Vanesa Quiroz] Coordinadora del proyecto. Ex coordinadora de cheLA (2004--2007), período en que desarrolló \textbf{EduSol} (programa de educación con Software Libre para escuelas de Parque Patricios) e integró el programa de investigación \textbf{sur@norte@sur} (Latin American Digital Cultures, UCLA). 20 años de trayectoria en tecnología comunitaria y cultura en América Latina. Radicada en Tres Cruces, Quebrada de Humahuaca. Fue docente de Tecnología en la \textbf{Tecnicatura de Desarrollo Indígena del COAJ} (Centro de Organizaciones Aborígenes de Jujuy) y docente de diseño gráfico y video en la \textbf{Tecnicatura de Comunicación Social del IES de Humahuaca}, donde formó a muchos de los referentes comunitarios actuales de los parajes involucrados en el proyecto. Acompañó proyectos productivos en Azul Pampa en el marco de la carrera de Técnico Universitario en Valor Agregado en Proyecto Productivo. Responsable del diseño, gestión y hosting de los proyectos productivos y turísticos familiares y comunitarios que la red habilita, incluyendo identidad digital, integración WhatsApp--MQTT y sistemas de comunicación con el exterior.

  \item[Fabian Wagmister] Asesor estratégico. Director General de cheLA. Fundador de UCLA REMAP. Especialista en \textit{Cultural Civic Computing} y tecnología para el empoderamiento cívico en comunidades latinoamericanas.

  \item[Negra Xue] Coordinadora del nodo BioCulturas de cheLA. Enlace institucional entre el proyecto y la red cheLA/Fundación exACTa.

  \item[Ing. Edmundo Ramos] Asesor técnico en movilidad. Creador de la tecnología autoabasura.com. Responsable de la instalación y transferencia de conocimiento del gasificador de biomasa.

  \item[Técnico Meshtastic Argentina] Por definir. Responsable de la capacitación técnica en instalación, mantenimiento de nodos LoRa y configuración de gateways MQTT.

  \item[Instalador local] Por definir. Técnico con experiencia en trabajo en altura. Responsable del despliegue de nodos en campo.
\end{description}

% ═══════════════════════════════════════════════════════════════════════════════
\section{Información de Contacto}
% ═══════════════════════════════════════════════════════════════════════════════

\begin{tabular}{ll}
\textbf{Organización:}     & Colectivo Nodo Norte \\
\textbf{Entidad legal:}    & Fundación exACTa / cheLA \\
\textbf{Responsable:}      & Patricia Vanesa Quiroz \\
\textbf{Email:}            & \href{mailto:vaneq@colectivonodonorte.ar}{vaneq@colectivonodonorte.ar} \\
\textbf{Web org.:}         & \href{https://colectivonodonorte.ar}{colectivonodonorte.ar} \\
\textbf{Web institucional:}& \href{https://chela.org.ar}{chela.org.ar} \\
\textbf{Dirección:}        & Tres Cruces, Dpto. Humahuaca, Provincia de Jujuy, Argentina \\
\textbf{GitHub:}           & \href{https://github.com/Colectivo-Nodo-Norte}{github.com/Colectivo-Nodo-Norte} \\
\end{tabular}

\vfill
{\color{primary}\rule{\linewidth}{1pt}}
{\small\color{darkgray}\centering
  Documento generado en \LaTeX{} --- Versión 1.1 --- Mayo 2026\\
  Repositorio: \href{https://github.com/Colectivo-Nodo-Norte/Red-Yerba-Buena}{github.com/Colectivo-Nodo-Norte/Red-Yerba-Buena}
}

\end{document}
